\clearpage
\fancyhead[L]{\small CONNAISSANCE DES TEMS.}
\hypertarget{native-643-lower}{}
\begin{center}\small Göttingische gelehrte Anzeigen. No. 80. pp. 793..796. 17 May 1824.\\[3pt]\rule{.72\linewidth}{.3pt}\end{center}

\emph{Connaissance des tems ou des mouvemens célestes à l'usage des astronomes et des navigateurs pour l'an \textup{1826}. Publiée par le bureau des longitudes.} Paris 1823. Published by Bachelier. The standing sections, 216 pages, and the supplements, 110 pages. Large octavo.

The supplements to this volume consist of the following articles. An overview of the astronomical observations made at the Marseille Observatory in the years 1820..1822 by \textsc{Gambart} the younger. The Marseille Observatory has so far been equipped with few instruments of middling quality, so Mr.\ \textsc{Gambart} has rightly limited himself to observations suited to these circumstances. The almost invariably clear sky of Marseille is particularly inviting for the search for new comets. The good fortune that once so favoured \textsc{Pons}, who became famous through these searches, seems to remain with \textsc{Gambart} as well. Of the three comets of 1822, he was the very first to discover one and discovered another himself, at least without knowing of its earlier discovery at Marlia. \textsc{Gambart}'s valuable observations of these two comets are printed here in as much detail as one could wish. \textsc{Gambart} also made two observations of the second comet of 1822. They are all the more valuable because the only other person to observe this comet was \textsc{Caturegli} in Bologna. The remaining observations are of eclipses of Jupiter's satellites, occultations of stars, and the lunar eclipse of 2 August 1822. --- On the distribution of heat in a homogeneous ring of uniform thickness situated in a region whose temperature differs from point to point [by \textsc{Poisson}]. The subtle mathematical treatment of this fine paper makes it of interest chiefly to the geometer. It also has a bearing on practical astronomy, since, during observations, circle instruments are often subject to conditions such as those on which the paper \SourceLeaf{644} is based, and it is desirable to gain a precise idea of the resulting deformation of the instruments. --- On the speed of sound [by \textsc{Poisson}]. This is an exposition as thorough as it is lucid of the theory of this subject, first brought to completion by \textsc{Laplace}. Using the ratio of the specific heat of air at constant pressure to its specific heat at constant volume, as obtained from the interesting experiment by \textsc{Clement} and \textsc{Delorme}, brings the theory of the speed of sound into almost perfect agreement with empirical results. --- Observations of the first and third comets of 1822 at the Paris Observatory. The latter had been discovered by Mr.\ \textsc{Bouvard} himself, though a few days later than by \textsc{Pons} and \textsc{Gambart}. --- Results from observations with invariable pendulums, calculated by Mr.\ \textsc{Mathieu}. The naval lieutenant \textsc{Duperrey}, who made a voyage around the world aboard the corvette Lacoquille in 1822 and 1823, took two invariable pendulums with him. Their oscillations were observed in Paris before his departure and by Mr.\ \textsc{Duperrey} during a landing in the Malouines. Since the southerly latitude of the latter is only $2^\circ\,41'$ greater than the northerly latitude of Paris, all that can really be concluded from these experiments, which in other respects are discussed with great care, is that they are consistent with a flattening of $\frac{1}{308}$. Of greater importance are the experiments with an invariable pendulum, also calculated here, that \textsc{Brisbane} and \textsc{Rümker} carried out first in London and then in Paramatta in New Holland. A flattening of $\frac{1}{206}$ is calculated here from these experiments. --- On the effect of the Moon on the atmosphere, by Mr.\ \textsc{de Laplace}. Since 1815 the height of the barometer at the Paris Observatory has been observed four times daily, at 9 a.m., noon, 3 p.m. and 9 p.m. The daily variation in the barometer reading stands out so clearly from the random fluctuations that the difference between the readings at 9 a.m. and 3 p.m., when averaged for each individual month over six years of observations, never once changed sign (the difference averages 0.801 millimetres, with the morning reading being the higher). This prompted an attempt to determine the influence of the Moon on the barometer reading from a suitable selection of these observations. For this purpose, the observations near syzygies were compared with those near quadratures. The comparison yielded a small periodic inequality with a period of half a lunar day. Its maximum is only 0.0272 millimetres and occurs on the day of syzygy at $3^{\rm u}18^{\rm m}36^{\rm s}$. Some inaccuracies appear to have crept into the numerical part of this investigation. This and the interest of the subject prompted the reviewer to repeat the calculations, in part more precisely and in part by other methods. The result, however, is almost the same as that obtained by \textsc{Laplace}. The reviewer finds a maximum of 0.0269 millimetres, a time of $3^{\rm u}11^{\rm m}49^{\rm s}$ for that maximum, and a mean error of 0.0264 millimetres to be expected in the first determination. --- This volume concludes with a short paper by Mr.\ \textsc{Puissant} on calculating the azimuths measured with a repeating theodolite. The procedure proposed here is already known and is essentially identical to the one that Mr.\ \textsc{Soldner} discussed in the Denkschriften der Münchner Akademie for 1813. Our journal for 1815, pp.\ 449 ff.\ [p.\ 587 of this volume], may also be consulted on this point.
